\section{问题描述}

本项目关注的是\textbf{多机械臂搬运任务的运筹学调度求解与二臂/三臂运行模式选择}，不是机器人运动学、动力学建模或轨迹控制问题；机械臂在本文中主要作为\textbf{可分配的作业资源}参与调度优化。

因此，本文并不研究机械臂关节角、末端轨迹、动力学方程或控制器设计，而是将多机械臂搬运过程抽象为一个具有\textbf{任务分配、执行排序、资源互斥、安全区占用和多目标优化}特征的调度问题。本章主要回答三个问题：第一，实际搬运任务如何抽象为任务集合与执行模式；第二，时间、能耗和安全约束如何用数学符号表达；第三，如何建立同时考虑\textbf{最大完工时间、总能耗和负载均衡}的优化模型。

\subsection{问题场景}

考虑一个圆形作业区域，区域中随机放置若干待搬运方块。每个方块具有类型、重量、初始位置和目标收集盒。系统中可配置两只或三只机械臂，机械臂从各自的待机位置出发，将方块搬运到对应类型的收集盒中。四类方块的任务属性如下：

\begin{table}[H]
    \centering
    \caption{方块类型与任务属性}
    \begin{tabular}{ccccc}
        \toprule
        方块类型 & 任务类型 & 所需机械臂数 & 目标收集盒 & 相对搬运难度 \\
        \midrule
        Type1 & 单臂任务 & 1 & Box1 & 小而轻 \\
        Type2 & 单臂任务 & 1 & Box2 & 小但较重 \\
        Type3 & 双臂协同任务 & 2 & Box3 & 大且需协同 \\
        Type4 & 双臂协同任务 & 2 & Box4 & 大且较重，协同代价高 \\
        \bottomrule
    \end{tabular}
    \label{tab:block_type_task}
\end{table}

图 \ref{fig:workspace_example} 给出同一组任务实例在两臂系统和三臂系统下的工作空间示意。图中蓝色圆形区域表示方块随机生成的作业区域，橙色虚线圆表示中心安全参考区，方形标记表示机械臂基座，菱形标记表示四类收集盒，圆点表示待搬运方块，虚线表示方块到对应收集盒的搬运方向。方块旁标注的机械臂组合为该案例求解得到的任务分配结果。

\begin{figure}[H]
    \centering
    \begin{minipage}{0.48\textwidth}
        \centering
        \includegraphics[width=\linewidth]{workspace_two_arm_example.jpg}
        \par\small (a) 两臂系统示意
    \end{minipage}
    \hfill
    \begin{minipage}{0.48\textwidth}
        \centering
        \includegraphics[width=\linewidth]{workspace_three_arm_example.jpg}
        \par\small (b) 三臂系统示意
    \end{minipage}
    \caption{多机械臂分类搬运任务的工作空间示意}
    \label{fig:workspace_example}
\end{figure}

在两臂系统中，双臂协同任务只能由两只机械臂共同完成；在三臂系统中，双臂协同任务可以从三只机械臂中选择任意两只共同完成。因此，三臂系统一方面增加了\textbf{可选执行模式和并行能力}，可能缩短最大完工时间；另一方面也会带来额外的系统固定能耗和协同协调代价。本文后续案例分析正是围绕这一\textbf{时间--能耗权衡}展开。

此外，作业区域中心附近设置安全区。当某个任务的搬运路径经过中心安全区时，该任务组需要占用中心区资源。这里的中心区不是单只机械臂层面的互斥资源，而是任务组层面的共享安全资源：同一个双臂协同任务内部的两只机械臂允许同时进入中心区，但不同任务组不能同时占用中心区。

\subsection{集合、参数与执行模式}

设任务集合为
\[
    \mathcal{J}=\{1,2,\ldots,n\}.
\]
两臂系统和三臂系统的机械臂集合分别为
\[
    \mathcal{A}^{(2)}=\{\mathrm{arm1},\mathrm{arm2}\},\qquad
    \mathcal{A}^{(3)}=\{\mathrm{arm1},\mathrm{arm2},\mathrm{arm3}\}.
\]
下文若不特别区分两臂或三臂，统一记机械臂集合为 $\mathcal{A}$，其规模为
$|\mathcal{A}|=m$，其中 $m\in\{2,3\}$。

每个任务 $j\in\mathcal{J}$ 对应一个方块。为便于后续建模，表 \ref{tab:basic_symbols} 汇总主要任务参数。

\begin{table}[H]
    \centering
    \caption{任务与系统参数说明}
    \begin{tabular}{c >{\centering\arraybackslash}p{0.68\textwidth}}
        \toprule
        符号 & 含义 \\
        \midrule
        $p_j$ & 任务 $j$ 对应方块的初始位置 \\
        $q_j$ & 任务 $j$ 对应目标收集盒的位置 \\
        $b_j$ & 方块类型，$b_j\in\{1,2,3,4\}$ \\
        $r_j$ & 完成任务 $j$ 所需机械臂数量 \\
        $w_j$ & 方块重量 \\
        $d_j$ & 方块从初始位置到目标收集盒的搬运距离 \\
        \bottomrule
    \end{tabular}
    \label{tab:basic_symbols}
\end{table}

不同方块类型决定任务所需机械臂数量：
\[
    r_j=
    \begin{cases}
    1, & b_j\in\{1,2\},\\
    2, & b_j\in\{3,4\}.
    \end{cases}
\]
搬运距离定义为
\[
    d_j=\|p_j-q_j\|_2.
\]

对每个任务 $j$，根据可达性和所需机械臂数构造可行执行模式集合 $\mathcal{M}_j$。一个执行模式 $k\in\mathcal{M}_j$ 包含以下信息：
\[
    k=(j,A_k,t_k,e_k,z_k),
\]
其含义如表 \ref{tab:mode_symbol} 所示。

\begin{table}[H]
    \centering
    \caption{执行模式 $k$ 的组成}
    \begin{tabular}{c >{\centering\arraybackslash}p{0.68\textwidth}}
        \toprule
        符号 & 含义 \\
        \midrule
        $j$ & 该模式对应的任务编号 \\
        $A_k$ & 执行该模式占用的机械臂集合，满足 $A_k\subseteq\mathcal{A}$ 且 $|A_k|=r_j$ \\
        $t_k$ & 该模式下的任务处理时间 \\
        $e_k$ & 该模式下的任务处理能耗 \\
        $z_k$ & 是否经过中心安全区，$z_k=1$ 表示经过，$z_k=0$ 表示不经过 \\
        \bottomrule
    \end{tabular}
    \label{tab:mode_symbol}
\end{table}

\subsubsection{任务处理时间}

本文将机械臂完成一个任务的时间分为两部分：一是执行当前任务本身所需的处理时间，二是任务之间的空载转移时间。这里先定义任务处理时间；空载转移时间将在后续序列相关约束中单独建模。

任务处理时间只描述机械臂到达方块位置后，将方块搬运到目标收集盒所需的时间。它由负载搬运时间、抓取/放置服务时间和安全缓冲时间组成。设单臂负载速度为 $v_s$，双臂协同负载速度为 $v_d$，则
\[
    t_k=
    \begin{cases}
    \left(d_j/v_s+T_s+T_{\mathrm{safe}}\right)\cdot \kappa_t, & |A_k|=1,\\
    \left(d_j/v_d+T_d+T_{\mathrm{safe}}\right)\cdot \kappa_t, & |A_k|=2,
    \end{cases}
\]
其中 $d_j/v_s$ 或 $d_j/v_d$ 表示带负载搬运阶段的行走时间；$T_s,T_d$ 分别表示单臂和双臂任务的抓取、夹持稳定与放置服务时间；$T_{\mathrm{safe}}$ 为调度层面的等效安全缓冲时间，用于表示对准、同步、放置稳定等未单独细分的短时间余量。$\kappa_t$ 为时间离散化尺度，代码实现中取 $\kappa_t=100$，即将 1 秒连续时间转换为 100 个整数 tick，以便整数规划和分枝定界算法处理。

\subsubsection{任务处理能耗}

任务处理能耗采用等效能耗模型描述，用于在调度层面区分不同方块类型的搬运难度。该能耗并不是对电机电流、关节力矩等底层物理量的精确预测，而是根据方块重量、搬运距离、任务持续时间和双臂协同复杂度构造的综合评价指标。这里定义的处理能耗记为 $e_k$，即后文系统总能耗公式中 $\sum e_kx_k$ 对应的任务执行能耗项。

对任务 $j$ 的某个执行模式 $k$，其处理能耗可表示为
\[
    e_k=\kappa_e\left(
    \phi_j E_{jk}^{\mathrm{load}}
    +E_{jk}^{\mathrm{grasp}}
    +|A_k|h_j\frac{t_k}{\kappa_t}
    +\sigma_j\eta_m
    \right).
\]
其中，$E_{jk}^{\mathrm{load}}$ 为基础负载搬运能耗，与方块重量和搬运距离有关；$E_{jk}^{\mathrm{grasp}}$ 为抓取、夹持建立和放置动作产生的服务能耗；$|A_k|h_jt_k/\kappa_t$ 表示任务执行期间参与机械臂维持夹持与稳定搬运所需的保持能耗；$\sigma_j\eta_m$ 表示双臂协同任务中的额外同步协调代价及系统规模修正。$\kappa_e$ 为能耗离散化尺度，用于将连续能耗转换为整数能耗值，以便与整数规划和分支定界算法统一。

为刻画不同方块类型对能耗的影响，引入表 \ref{tab:energy_type_params} 中的类型相关参数。

\begin{table}[H]
    \centering
    \caption{方块类型相关能耗参数}
    \begin{tabular}{
        >{\centering\arraybackslash}p{0.18\textwidth}
        >{\centering\arraybackslash}p{0.24\textwidth}
        >{\centering\arraybackslash}p{0.46\textwidth}
    }
        \toprule
        参数 & 名称 & 含义 \\
        \midrule
        $\phi_j$ & 载荷修正系数 & 放大搬运难度对能耗的影响 \\
        $h_j$ & 保持功率 & 单位时间内的夹持与稳定搬运代价 \\
        $\sigma_j$ & 同步能耗 & 双臂协同任务的一次性协调代价 \\
        \bottomrule
    \end{tabular}
    \label{tab:energy_type_params}
\end{table}

对于 Type1 和 Type2 单臂任务，任务不涉及双臂同步，因此令 $\sigma_j=0$，且 $|A_k|=1$。对于 Type3 和 Type4 双臂协同任务，$\sigma_j$ 取正值，$|A_k|=2$，用于刻画两只机械臂在同步支撑、同步移动和协调放置过程中的额外能耗。两臂系统中 $\eta_m=1$；三臂系统虽然每个双臂任务仍只使用两只机械臂，但可选组合更多、调度协调更复杂，因此 $\eta_m$ 可取略大于 1 的值。其中 Type4 方块的搬运难度更高，因此其 $\phi_j$、$h_j$ 或 $\sigma_j$ 可设置为更大的取值。

\subsection{决策变量}

为了同时描述任务模式选择、时间安排和机械臂作业顺序，定义如下决策变量。

\begin{itemize}
    \item $x_k\in\{0,1\}$：若任务执行模式 $k$ 被选择，则 $x_k=1$，否则为 0。
    \item $S_k,E_k$：模式 $k$ 对应任务处理段的开始时间和结束时间。
    \item $y_{a,u,v}\in\{0,1\}$：对机械臂 $a$，若其作业顺序中模式 $u$ 后紧接模式 $v$，则 $y_{a,u,v}=1$。虚拟模式 $0$ 表示机械臂的待机位置，因此 $y_{a,0,v}=1$ 表示 $v$ 是机械臂 $a$ 的首个任务，$y_{a,u,0}=1$ 表示 $u$ 是其末个任务。
    \item $C_{\max}$：所有任务的最大完工时间。
    \item $E_{\mathrm{tot}}$：系统总能耗。
    \item $L_a$：机械臂 $a$ 的总负载时间，包括任务处理时间和空载转移时间。
    \item $B$：机械臂负载不均衡度，即最大负载与最小负载之差。
\end{itemize}

其中，$y_{a,u,v}$ 用于描述序列相关 setup。机械臂执行完一个任务后，需要从上一任务目标收集盒移动到下一任务的方块初始位置；若该任务是该机械臂的第一个任务，则需要从机械臂待机位置移动到该方块位置。这部分空载转移时间和能耗都与任务顺序有关。为简化后续公式，定义参数
\[
    \delta_{a,k}=\begin{cases}
    1, & a\in A_k,\\
    0, & a\notin A_k,
    \end{cases}
    \qquad
    \mathcal M_a=\{k:\delta_{a,k}=1\}.
\]

\subsection{约束条件}
\begin{enumerate}
\item \textbf{任务唯一执行模式约束}

每个任务必须且只能选择一种可行执行模式：
\[
    \sum_{k\in\mathcal{M}_j}x_k=1,\quad \forall j\in\mathcal{J}.
\]
该约束保证每个方块都被搬运一次，且不会被多个模式重复执行。

\item \textbf{时间一致性约束}

若模式 $k$ 被选择，则其结束时间等于开始时间加处理时间：
\[
    E_k=S_k+t_k,\quad \forall k \text{ with } x_k=1.
\]
在整数规划实现中，未被选择的模式通过可选区间变量处理；在自编算法中，未被选择的模式不会进入调度序列。

\item \textbf{机械臂资源互斥约束}

同一只机械臂在同一时刻不能执行两个任务。若两个被选择模式 $u,v$ 均占用机械臂 $a$，则必须满足二者在时间轴上不重叠：
\[
    [S_u,E_u)\cap [S_v,E_v)=\varnothing,\quad
    \forall a\in A_u\cap A_v.
\]
基础模型中该约束通过 NoOverlap 资源约束实现；自编求解器中则通过维护每只机械臂的可用时间递推生成可行调度。

\item \textbf{序列相关空载转移约束}

首先需要使顺序变量与模式选择变量一致。对每只机械臂，任一被选模式恰有一个前驱和一个后继：
\[
    \sum_{\substack{u\in\mathcal M_a\cup\{0\}\\u\ne v}}y_{a,u,v}
    =\delta_{a,v}x_v,
    \qquad \forall a\in\mathcal A,\ v\in\mathcal M_a,
\]
\[
    \sum_{\substack{v\in\mathcal M_a\cup\{0\}\\v\ne u}}y_{a,u,v}
    =\delta_{a,u}x_u,
    \qquad \forall a\in\mathcal A,\ u\in\mathcal M_a.
\]
同时令虚拟节点 $0$ 具有至多一条出弧和一条入弧，并施加单回路约束以消除不经过虚拟节点的子回路。代码中上述关系由 \texttt{AddCircuit} 统一实现；当机械臂未承担任务时，采用虚拟自环 $(0,0)$。

设机械臂 $a$ 执行完模式 $u$ 后紧接执行模式 $v$，则模式 $v$ 的开始时间必须不早于模式 $u$ 的结束时间加空载转移时间：
\[
    S_v\ge E_u+\tau_{u,v},\quad \text{if } y_{a,u,v}=1,
    \qquad u,v\in\mathcal M_a.
\]
其中
\[
    \tau_{u,v}=\frac{\|q_u-p_v\|_2}{v_0}\cdot \kappa_t,
\]
$v_0$ 为空载移动速度。若 $v$ 是机械臂 $a$ 的第一个任务，则
\[
    S_v\ge \tau^a_{0,v}
    =\frac{\|o_a-p_v\|_2}{v_0}\cdot \kappa_t,
    \quad \text{if }y_{a,0,v}=1,
\]
其中 $o_a$ 为机械臂 $a$ 的待机位置。

相应的空载转移能耗记为
\[
    \epsilon_{u,v}=\alpha_0\|q_u-p_v\|_2\cdot \kappa_e,
    \qquad
    \epsilon^a_{0,v}=\alpha_0\|o_a-p_v\|_2\cdot \kappa_e,
\]
首任务从待机位置出发时同理计算。该项将计入系统总能耗。

\item 中心安全区互斥约束

参数 $z_k\in\{0,1\}$ 表示模式 $k$ 的负载搬运路径是否经过中心安全区。对于两个不同任务的已选模式 $u,v$，若 $z_u=z_v=1$，则其中心区占用时间不能重叠：
\[
    [S_u,E_u)\cap [S_v,E_v)=\varnothing,
    \quad \text{if }x_u=x_v=z_u=z_v=1,
    \quad j(u)\ne j(v).
\]
该约束体现中心安全区在任务组层面的互斥。需要强调的是，对于同一个双臂协同任务，其内部两只机械臂属于同一任务组，不构成彼此冲突。

\item 最大完工时间、总能耗与负载均衡约束

最大完工时间满足
\[
    C_{\max}\ge E_k,\quad \forall k \text{ with } x_k=1.
\]

系统总能耗由处理能耗、空载转移能耗和系统级固定开销组成：
\[
    E_{\mathrm{tot}}
    =
    \sum_{j\in\mathcal{J}}\sum_{k\in\mathcal{M}_j}e_kx_k
    +\sum_{a\in\mathcal{A}}\sum_{\substack{u,v\in\mathcal M_a\\u\ne v}}
    \epsilon_{u,v}y_{a,u,v}
    +\sum_{a\in\mathcal A}\sum_{v\in\mathcal M_a}
    \epsilon^a_{0,v}y_{a,0,v}
    +mE_{\mathrm{sys}},
\]
其中 $E_{\mathrm{sys}}$ 为单只机械臂的启动、上电、伺服保持和控制通信等效固定能耗。

机械臂负载 $L_a$ 包括该机械臂承担的处理时间和空载转移时间：
\[
    L_a=
    \sum_k\delta_{a,k}t_kx_k
    +\sum_{\substack{u,v\in\mathcal M_a\\u\ne v}}\tau_{u,v}y_{a,u,v}
    +\sum_{v\in\mathcal M_a}\tau^a_{0,v}y_{a,0,v}.
\]
负载均衡指标定义为
\[
    B=\max_{a\in\mathcal{A}}L_a-\min_{a\in\mathcal{A}}L_a.
\]
\end{enumerate}

\subsection{目标规划模型}

本问题同时关注完成效率、运行能耗和机械臂利用均衡，三个目标按时间、能耗、负载均衡的顺序具有严格优先级。统一目标框架为
\[
    \operatorname{lexmin}\bigl(C_{\max},E_{\mathrm{tot}},B\bigr).
\]
其具体含义为：
\[
    P_1:\quad \min C_{\max},
\]
\[
    P_2:\quad \min E_{\mathrm{tot}},
    \qquad \text{s.t. }C_{\max}=C_{\max}^{\star},
\]
\[
    P_3:\quad \min B,
    \qquad \text{s.t. }C_{\max}=C_{\max}^{\star},\quad
    E_{\mathrm{tot}}=E_{\mathrm{tot}}^{\star}.
\]
即先保证整体任务尽快完成，再在不牺牲时间目标的前提下降低能耗，最后改善机械臂负载均衡。各求解方法遵循相同的优先级，但实现形式不同。

\subsubsection{CP-SAT 的偏差变量目标规划形式}

基础 CP-SAT 方法采用带偏差变量的目标规划。设 $G_T,G_E,G_B$ 为三个理想目标值，引入非负偏差变量 $d_i^+,d_i^-$：
\[
\begin{aligned}
    C_{\max}+d_1^- -d_1^+ &=G_T,\\
    E_{\mathrm{tot}}+d_2^- -d_2^+ &=G_E,\\
    B+d_3^- -d_3^+ &=G_B.
\end{aligned}
\]
其中 $d_i^+$ 表示超过理想目标的正偏差。令 $\underline{T}_j$ 为任务 $j$ 的最短处理时间与最小初始接近时间之和，$\underline e_j$ 和 $\underline\epsilon_{0,j}$ 分别为最低处理能耗和最小初始接近能耗，则程序自动估计
\[
\begin{aligned}
    G_T&=\max\left\{\max_j\underline T_j,
    \left\lceil\frac{\sum_j\underline T_j}{m}\right\rceil\right\},\\
    G_E&=\sum_j\underline e_j+\sum_j\underline\epsilon_{0,j}+mE_{\mathrm{sys}},
    \qquad G_B=0.
\end{aligned}
\]
其中 $G_B=0$ 表示理想负载完全均衡。取 $M=10^5$，CP-SAT 分三阶段求解：
\[
    P_1:\quad \min\; Md_1^+ +C_{\max}.
\]
得到 $d_1^{+\star}$ 和 $C_{\max}^{\star}$ 后固定二者，再求解
\[
    d_1^+=d_1^{+\star},\qquad
    C_{\max}=C_{\max}^{\star},
    \qquad
    P_2:\quad \min\; Md_2^+ +E_{\mathrm{tot}}.
\]
第三阶段保持时间最优，并允许第二阶段能耗在很小的工程容差内波动：
\[
    d_2^+\le d_2^{+\star}+\Delta_d,
    \qquad
    E_{\mathrm{tot}}\le E_{\mathrm{tot}}^{\star}+\Delta_E,
    \qquad
    P_3:\quad \min\; Md_3^+ +B.
\]
代码中 $\Delta_d=\Delta_E=10$；若后续阶段在时限内未返回可用解，则回退到上一阶段结果。

\subsubsection{自编算法的字典序形式}

融合剪枝的隐枚举法、启发式算法和热启动后的隐枚举法不显式建立偏差变量，而是直接按顺序比较
\[
    \bigl(C_{\max},E_{\mathrm{tot}},B\bigr),
\]
进行字典序更新。因此，两类实现的数学形式不同，但目标优先级一致。

综上，本文的核心模型可以理解为一个带有可选执行模式、并行机器资源约束、序列相关 setup、中心安全区互斥和多目标优先级的调度优化模型。后续第 3 章将在该模型基础上，分别设计 CP-SAT求解器、融合剪枝的隐枚举法、启发式算法和热启动后的隐枚举法四类求解方法，并在第 4 章进一步讨论速度--能耗层的非线性扩展与二臂/三臂模式选择结果。
